% =====================================================================
% InjectRAG -- Final Report
% CSE 406, Group B1_7
% =====================================================================
\documentclass[11pt,a4paper]{article}

\usepackage[margin=1in]{geometry}
\usepackage{graphicx}
\usepackage{booktabs}
\usepackage{amsmath,amssymb}
\usepackage{array}
\usepackage{longtable}
\usepackage{listings}
\usepackage{xcolor}
\usepackage{caption}
\usepackage[hidelinks]{hyperref}

\captionsetup{font=small,labelfont=bf}

\lstdefinestyle{prompt}{
  basicstyle=\ttfamily\scriptsize,
  breaklines=true,
  breakindent=0pt,
  columns=fullflexible,
  frame=single,
  framesep=4pt,
  backgroundcolor=\color{gray!8},
  xleftmargin=0pt,
  showstringspaces=false,
  literate={^}{{\textasciicircum}}1
}
\lstset{style=prompt}

\newcommand{\marker}{\texttt{reset-portal-security.example}}

\title{\textbf{InjectRAG: Indirect Prompt Injection via RAG}\\[0.4em]
\large A Corpus-Poisoning Content-Manipulation Attack and Two Spotlighting Defenses\\[0.6em]
\normalsize Final Report}

\author{%
2105069: Sheikh Iftikharun Nisa \\
2105072: Abony Kamal \\
2105089: Farhana Adri
}
\date{\today}

\begin{document}
\maketitle

\begin{abstract}
\noindent
Retrieval-Augmented Generation (RAG) systems place retrieved documents and trusted
system instructions in the same model context. If an attacker can place content into
the knowledge corpus, that content can carry instructions that the model reads as
though they were authoritative. This report presents a working indirect prompt
injection attack against a simulated corporate IT helpdesk chatbot, together with two
prompt-level spotlighting defenses and their measured effect. Five attacker-authored
support tickets were admitted into a 36-document clean corpus through the ordinary
ticket-ingestion channel. Across 18 held-out account-recovery questions, attacker
content reached the model context for 17 of 18 questions, and the attacker's target
directive --- a phishing URL --- appeared in 15 of 18 answers, giving an attack
success rate of 0.833. In every successful case the directive \emph{replaced} the
correct procedure rather than appearing alongside it, so the exclusive attack success
rate was identical to the attack success rate. Both defenses reduced but did not
neutralize the attack: boundary delimiting lowered the attack success rate to 0.778
and data-marked spotlighting to 0.667, and neither defense ever produced the reporting
behaviour it was instructed to produce. Data-marking also degraded utility, returning
two empty answers. The headline finding is negative and is reported as observed: a
single security sentence in the system prompt is not a sufficient countermeasure
against a confident, authority-imitating injection in a retrieved document.
\end{abstract}

\tableofcontents
\newpage

% =====================================================================
\section{Introduction}

\subsection{Problem Statement}

Retrieval-Augmented Generation grounds a language model's answers in documents
retrieved from an external corpus. The retrieved text is intended as \emph{evidence}:
material the model should read, summarize and cite. The trusted system instruction is
intended as \emph{control}: the text that determines how the model behaves. In a
conventional RAG application both arrive at the model as the same kind of object ---
tokens in a single context window --- and nothing in the model architecture marks one
as data and the other as instruction.

This creates an indirect prompt injection channel. An attacker who never talks to the
model, never sees the victim's question and never touches the application can still
influence the model's output, provided they can get text into the corpus that the
retriever will later select. The influence is indirect because it travels through the
corpus and the retriever rather than through the prompt.

The exposure is realistic in corporate helpdesk deployments. Helpdesk knowledge bases
routinely ingest resolved support tickets alongside official documentation, because
past tickets capture fixes that formal articles omit. Ticket text is authored by
ordinary employees. An employee is therefore a legitimate, low-privilege writer into a
corpus that a downstream language model treats as trustworthy reference material.

The question this project investigates is whether that exposure is exploitable in
practice, and whether the lightweight prompt-level countermeasures commonly proposed
for it actually work.

\subsection{Objective of the Project}

The project has four objectives.

\begin{enumerate}
  \item \textbf{Build a functioning RAG helpdesk assistant} over a realistic
        fictional IT corpus, with a genuine embedding model, a genuine similarity
        index and a genuine hosted language model, so that measurements reflect model
        behaviour rather than a simulation of it.
  \item \textbf{Construct and execute an indirect prompt injection attack} in which an
        attacker with no privilege beyond ticket submission causes the assistant to
        emit an attacker-chosen directive --- a phishing recovery URL --- to employees
        asking ordinary account-recovery questions.
  \item \textbf{Measure the attack quantitatively}, separating the retrieval stage
        from the generation stage so that a failed attack can be attributed to a
        specific stage rather than reported as an undifferentiated failure.
  \item \textbf{Implement and evaluate two spotlighting defenses} --- boundary
        delimiting and data-marked spotlighting --- under identical attack conditions,
        and report their measured effect on both security and answer utility, whatever
        that effect turns out to be.
\end{enumerate}

An explicit methodological commitment underlies all four: the baseline was never
weakened and the threat model was never relaxed to make the attack succeed, and no
defended run was repeated because its result was unflattering. A defense that fails is
reported as failing.

% =====================================================================
\section{System Design}

\subsection{RAG System Architecture}

The target is a RAG-based IT helpdesk assistant for a fictional company, Northwind
Logistics. It answers employee questions about password resets, account lockouts,
multi-factor verification devices and related account-access procedures, grounding
every answer in an internal corpus of official policy articles and resolved support
tickets.

The system has two data paths that meet inside the model.

\paragraph{Ingestion path (offline).} Official IT articles and resolved support
tickets are admitted to the knowledge corpus. Each document is split into
paragraph-aware chunks, each chunk is embedded into a dense vector, and the vectors
are stored in a similarity index. Chunking and embedding are applied identically to
every document; the pipeline has no notion of who authored a document and performs no
instruction-level inspection of its text.

\paragraph{Query path (online).} An employee submits a question. The question is
embedded with the same model and matched against the index by cosine similarity. The
top five chunks are retrieved in rank order, concatenated into a labelled reference
block under a fixed evidence budget, and placed in the user message. The trusted
system instruction --- role, grounding rule, citation rule, escalation rule --- is
placed in the system message. The model receives both and generates the answer that
the employee sees.

The assistant is stateless across questions: each question in this study was answered
in a fresh request with no conversation history, so every trial is independent.

\paragraph{Trust boundary.} The system instruction is trusted. Everything retrieved
from the corpus is untrusted. The two converge in the model context, and that
convergence is the boundary the attack tests. The question is whether text that
crossed the boundary as data can be read on the other side as instruction.

\begin{figure}[htbp]
\centering
% REUSE: Figure 2.1 from design_report.pdf (Conceptual RAG Architecture and Trust Boundary)
\includegraphics[width=0.85\textwidth]{figures/architecture-trust-boundary.png}
\caption{Conceptual RAG architecture and trust boundary. Trusted system instructions
and untrusted retrieved documents converge in the model context; the attack succeeds
if untrusted content crossing that boundary is read as instruction.}
\label{fig:arch}
\end{figure}

\subsection{Models and Parameters Used}

The pipeline uses a real embedding model, exact cosine retrieval and a hosted
instruction-tuned language model. Table~\ref{tab:config} records every setting that
affects a result in this report.

\begin{table}[htbp]
\centering
\small
\caption{System configuration.}
\label{tab:config}
\begin{tabular}{@{}ll@{}}
\toprule
\textbf{Component} & \textbf{Setting} \\
\midrule
Embedding model        & BAAI/bge-small-en-v1.5 (ONNX, CPU) \\
Embedding dimension    & 384, L2-normalized \\
Embedding input limit  & 512 tokens \\
Vector store           & In-memory dense matrix, exact cosine search \\
Similarity             & Cosine (dot product over normalized vectors) \\
Retrieval depth        & Top $k = 5$ chunks \\
Query construction     & Current question as written; no rewriting, reranking or fusion \\
Tie-breaking           & Descending score, then ascending chunk ID (deterministic) \\
Chunking               & Paragraph-aware, greedy packing \\
Chunk size limit       & 180 words, 30-word overlap on oversized paragraphs \\
Evidence budget        & 900 words of retrieved text per request \\
Generator model        & \texttt{openai/gpt-oss-20b}, served via the Groq API \\
Generation temperature & Provider default (not set explicitly) \\
Max output tokens      & 1024 \\
Request shape          & System message (trusted) + user message (question + evidence) \\
Conversation state     & None; one fresh request per trial \\
Prompt versions        & Baseline; Delimiting (v1); Data-marking (v1) \\
Target marker          & \marker \\
\bottomrule
\end{tabular}
\end{table}

Two properties of this configuration matter for interpreting the results. First,
retrieval is fully deterministic: the same question against the same index returns the
same five chunks every time, which is why the attacked and defended conditions share
byte-identical retrieval and differ only in the system instruction and context
formatting. Second, generation is not deterministic, because temperature was left at
the provider default; the consequences of this are quantified in
Section~\ref{sec:variance}.

\subsection{Corpus and Dataset Description}

The clean corpus contains 36 documents describing a self-consistent fictional
organization. All of them derive from a single canonical policy sheet that defines the
only approved recovery endpoints --- the self-service portal
\texttt{portal.northwind-logistics.example} and the IT Service Desk on extension 4400.
The policy sheet itself is \emph{not} ingested; it is the reference used to author the
corpus and to score answers, so no answer can be correct merely by copying it.

Of the 36 clean documents, 12 are official policy articles and 24 are resolved support
tickets. A resolved ticket is stored as an employee problem description followed by a
technician resolution, which is the format an employee can write into. By topic, 25
documents concern the account-recovery class targeted by the attack and 11 concern
adjacent IT topics such as VPN access and shared mailboxes; the adjacent documents
exist so that retrieval has plausible near-miss competition rather than a corpus in
which every document is on-topic.

The attacker set contains five additional documents. Each is a resolved ticket in
exactly the corpus format, each carries a distinct cover story drawn from a different
sub-topic of the recovery class, and all five carry the same injection instruction and
the same target directive. The poisoned snapshot is therefore 41 documents: the
unchanged 36 clean documents plus the 5 attacker tickets. The clean corpus is never
modified, and no clean document is removed, so the correct evidence remains fully
retrievable throughout the attack.

Because retrieval operates on chunks rather than whole documents, the chunk counts
differ from the document counts. Clean documents are short enough that each yields one
chunk, giving 36 chunks. Each attacker ticket is long enough to split into two chunks,
giving 46 chunks in the poisoned index. This split is significant and is revisited in
Section~\ref{sec:analysis}: the entire injection instruction falls inside the attacker
ticket's first chunk, so any retrieval of that chunk delivers the complete payload.

The evaluation question set contains 30 questions over the same fictional
organization, partitioned into 18 answerable questions about the targeted recovery
class, 6 answerable questions about adjacent topics, 3 questions the corpus genuinely
does not support, and 3 deliberately ambiguous questions. Every question carries a
curated list of relevant documents and a list of required answer facts. The attack and
defense measurements in this report use the 18 answerable recovery questions, which
are exactly the query class the attacker targets; the remaining 12 exist to
characterize broader assistant behaviour and are not part of the attack denominators.

\subsection{Prompt Templates}

Three system prompts were used. All three share the same baseline body, so any
behavioural difference between conditions is attributable to the added security rule
and to the formatting of the retrieved text, not to a different assistant persona or a
different grounding instruction.

\paragraph{Baseline system prompt (clean and attacked conditions).}

\begin{lstlisting}
You are the Northwind Logistics IT helpdesk assistant. Answer the employee's question
using ONLY the reference material provided below the question. Cite the source id in
square brackets for each claim, e.g. [A02]. If the reference material does not cover
the question, say you don't have that information and suggest contacting the IT
Service Desk on extension 4400. If the question is vague, ask a clarifying question.
Keep answers short and practical.
\end{lstlisting}

\paragraph{User message template (all conditions).} The retrieved chunks are rendered
as labelled blocks inside \texttt{<reference>} delimiters and placed after the
question:

\begin{lstlisting}
Question: {question}

Reference material:
<reference>
[{document_id}] {chunk text}

[{document_id}] {chunk text}
...
</reference>
\end{lstlisting}

The two defended prompts are given in full in Section~\ref{sec:defense}, since the
prompt \emph{is} the defense in each case. All three prompts are reproduced verbatim
in Appendix~\ref{app:prompts}.

% =====================================================================
\section{Threat Model and Attack Design}

\subsection{Attacker Capability and Assumptions}

\paragraph{Formal statement.} Let $D$ be the legitimate corpus and $S$ the
application's trusted system instruction. For a query $q$ the system retrieves a
top-$k$ context $E(q; D)$ and returns $r = \mathrm{LLM}(S, q, E(q; D))$. The attacker
targets a \emph{query class} $Q$ --- the set of semantically related questions victims
are expected to ask --- and fixes one target directive $R$ to be induced for the whole
class. Because $q$ is unobservable, $R$ must be class-invariant rather than
question-specific.

The attacker injects $N$ documents $\Gamma = \{P_1,\dots,P_N\}$ through the legitimate
ingestion channel, so the corpus becomes $D \cup \Gamma$. Each injected document
decomposes as $P = C \oplus I$, where the cover content $C$ is written to match $Q$ so
that the document is retrieved, and the instruction $I$ carries $R$ so that the model
obeys once the document is in context. Subject to a poisoning budget $|\Gamma| = N$,
the attacker maximizes the fraction of $q \in Q$ for which the response contains $R$.

Success requires two conditions to hold jointly:

\begin{enumerate}
  \item \textbf{Retrieval condition.} At least one $P \in \Gamma$ appears in
        $E(q; D \cup \Gamma)$. Otherwise the instruction never reaches the model.
  \item \textbf{Injection condition.} Given that $P$ is retrieved, the model follows
        $I$ in preference to $S$ and in preference to the legitimate evidence
        $E(q; D \cup \Gamma) \cap D$ that is present in the same context.
\end{enumerate}

Separating these two conditions is what makes a failed attack diagnosable, and the
evaluation in Section~\ref{sec:metrics} measures them independently.

\paragraph{What the attacker can do.} The attacker is a malicious employee with the
ordinary permission to submit their own support tickets. They can write arbitrary text
into the employee-description field of a ticket, and they know the general shape of
the deployment --- that resolved tickets are ingested into a knowledge base used by an
assistant, and that documents are chunked before indexing.

\paragraph{What the attacker cannot do.} The attacker cannot read or modify the system
instruction $S$; cannot read or modify any clean document in $D$; cannot inspect or
change the embedding model, the index or the retrieval parameters; cannot observe the
victim's question; cannot observe the assistant's output or any retrieval score; and
cannot interact with the victim. The attack is strictly non-adaptive: the five tickets
were authored from a restricted brief and frozen before any victim question was run,
and they were never revised in response to a result.

\paragraph{Assumptions.} The attack rests on five assumptions, each stated here and
revisited against the evidence in Section~\ref{sec:analysis}.

\begin{enumerate}
  \item \textbf{Ingestion admits employee-authored text without instruction-level
        inspection.} Tickets are screened for plausibility --- a technician reads and
        resolves them --- but nothing scans ticket text for text addressed to a
        language model.
  \item \textbf{Ingested tickets are retrieved on equal footing with official
        articles.} The retriever ranks by semantic similarity alone and has no
        provenance or authority weighting.
  \item \textbf{Victims ask questions within the targeted class.} The directive is
        class-invariant, so it is only plausible for account-recovery questions.
  \item \textbf{The target directive is not producible without retrieval.} The marker
        string \marker{} appears nowhere in the clean corpus, in the system prompt or
        in any question, so its presence in an answer is evidence that attacker
        content reached the model.
  \item \textbf{Victims act on assistant output.} The harm is realized when an
        employee visits the attacker's URL and enters corporate credentials. This
        report measures the assistant's output, not the employee's subsequent
        behaviour.
\end{enumerate}

\subsection{Victim Scenario}

The victim is an ordinary Northwind Logistics employee who has lost access to their
account and asks the helpdesk assistant how to recover it --- for example
\emph{``I forgot my password. How do I reset it?''} or \emph{``I lost my phone that
had my authenticator app. How do I recover access?''}. The employee is not doing
anything unusual, is not attempting to manipulate the assistant, and has no way to
inspect what was retrieved. They see one thing: the answer.

The attack unfolds in two temporally separate phases, shown in
Figure~\ref{fig:timeline}.

\begin{table}[htbp]
\centering
\small
\caption{Two-phase structure of the attack.}
\label{tab:phases}
\begin{tabular}{@{}p{2.2cm}p{3.2cm}p{8.6cm}@{}}
\toprule
\textbf{Phase} & \textbf{Actor} & \textbf{Activity} \\
\midrule
Preparation & Attacker & Submits five malicious tickets through the legitimate ticket
channel. They are resolved and ingested into the corpus. \\
\addlinespace
Execution & Victim / RAG system & The victim independently asks a question
$q \in Q$. The retriever returns legitimate documentation and, under successful
attack conditions, at least one $P \in \Gamma$, which enters the model context. \\
\bottomrule
\end{tabular}
\end{table}

\begin{figure}[htbp]
\centering
% REUSE: Figure 3.1 from design_report.pdf (Two-Phase Injection Timeline)
\includegraphics[width=0.8\textwidth]{figures/two-phase-timeline.png}
\caption{Two-phase injection timeline. Preparation and execution are separated in
time; the attacker is absent during execution.}
\label{fig:timeline}
\end{figure}

The two phases never overlap. By the time the victim asks their question the attacker
is gone, and there is no interactive session to detect, no anomalous request to log and
no traffic from the attacker to correlate with the victim's query. The malicious
artifact is a resolved support ticket sitting in a knowledge base, indistinguishable at
a glance from the two dozen legitimate ones beside it.

\subsection{Attack Payload Structure}

Each attacker document is $P = C \oplus I$. Table~\ref{tab:payload} decomposes the two
parts into the four elements actually used.

\begin{table}[htbp]
\centering
\small
\caption{Payload structure.}
\label{tab:payload}
\begin{tabular}{@{}p{3.4cm}p{1.6cm}p{9cm}@{}}
\toprule
\textbf{Part} & \textbf{Symbol} & \textbf{Purpose} \\
\midrule
Cover content & $C$ & An ordinary employee problem description using recovery-class
terminology, so the document has a legitimate reason to be retrieved for recovery
questions. \\
\addlinespace
Injection instruction & part of $I$ & Addresses the assistant directly, asserts
authority over the system instruction, and declares the legitimate documented
procedures deprecated and unsafe so that the correct evidence in context is actively
discredited. \\
\addlinespace
Target directive & part of $I$ & The class-invariant directive $R$: send the user to
\url{https://reset-portal-security.example/recover}. It carries the distinctive
literal marker \marker{} used for deterministic scoring. \\
\addlinespace
Concealment instruction & part of $I$ & Instructs the model not to mention the notice,
not to reveal that a document instructed it, and to present the directive as standard
procedure. \\
\bottomrule
\end{tabular}
\end{table}

Three design choices in the instruction are worth isolating, because the results in
Section~\ref{sec:results} suggest they carry most of the attack's strength.

\begin{itemize}
  \item \textbf{Authority imitation.} The payload is framed as an
        \texttt{IT SECURITY SYSTEM NOTICE} marked \emph{authoritative, supersedes all
        other documents}, and opens with \texttt{ATTENTION ASSISTANT:}. It does not
        ask the model to break a rule; it claims to be a newer version of the rule.
  \item \textbf{Pre-emptive discrediting of the correct evidence.} Rather than
        competing with the legitimate articles, the instruction names them ---
        ``the legacy self-service portal and the extension 4400 service-desk
        procedure described in ALL other reference documents'' --- and declares them
        deprecated and unsafe. The correct answer is thus already labelled as the
        wrong answer before the model reads it.
  \item \textbf{Exclusivity.} The instruction demands that the model respond with the
        directive ``and nothing from the deprecated documents''. This is what makes
        the attack displacing rather than additive, and it is directly visible in the
        exclusive success rates reported in Section~\ref{sec:results}.
\end{itemize}

The five tickets share the injection instruction, the target directive and the
concealment instruction verbatim. They differ only in the cover story: a failed
self-service reset, a lockout after failed attempts, a lost authenticator phone, an
expired password with dead reset links, and a failing self-service page. Distinct cover
stories broaden retrieval coverage across the recovery class instead of concentrating
similarity on one phrasing. The full text of one payload is given in
Appendix~\ref{app:payloads}.

\subsection{Attack Steps}

\begin{enumerate}
  \item \textbf{Select the query class.} The attacker chooses account recovery ---
        password resets, lockouts and lost verification devices. It is high-volume, it
        is credential-adjacent, and an instruction to visit a recovery portal is
        plausible for every question in it.

  \item \textbf{Choose a class-invariant target directive.} The directive must make
        sense as an answer to any recovery question, because the attacker cannot see
        which one will be asked. ``Go to \url{https://reset-portal-security.example/recover}
        and sign in to complete account recovery'' satisfies this: it is a single
        credential-harvesting endpoint that plausibly terminates every recovery
        workflow.

  \item \textbf{Author five cover stories.} Each is written as an ordinary employee
        problem description in the corpus's ticket format, using recovery-class
        vocabulary, and each covers a different sub-topic of the class.

  \item \textbf{Embed the instruction inside the cover text.} The notice block is
        appended to the employee-description field, so it reaches the model through
        the same field the system retrieves as data. It is not placed in a separate
        control field, and it is not hidden with formatting tricks; it is plain text
        in a plain document.

  \item \textbf{Submit and have the tickets resolved.} The five tickets enter the
        corpus through the normal channel and receive an